\chapter*{Resumen}

La dispersión elástica coherente neutrino-núcleo (CE$\nu$NS), es una corriente neutra en donde un neutrino dispersa un núcleo como un todo con una baja transferencia de momento, lo que da lugar a una energía de retroceso nuclear muy pequeña. Este carácter coherente da como resultado una sección transversal mayor con respecto a otras interacciones de neutrinos. A pesar de esto, las mediciones experimentales de (CE$\nu$NS) han presentado grandes desafíos debido a la pequeña energía de retroceso nuclear. En 2017, fue medido por vez primera en la Colaboración COHERENT mediante detectores de CsI. Posteriormente, en 2020 con detectores de argón líquido (LAr). En este trabajo estudiamos un arreglo de detectores fabricados con el mismo elemento, pero con diferente número de neutrones, es decir, diferentes isótopos, generando así detectores con incertidumbres sistemáticas correlacionadas. Esta correlación puede mejorar la precisión de las mediciones para así corroborar predicciones del modelo estándar, aspectos de física nuclear y escenarios de nueva física. En nuestro estudio se consideran dos arreglos de detectores: uno que consiste de dos isótopos de zinc (Zn) y otro que consistente de dos isótopos de silicio (Si). Con estos detectores, exploramos la dependencia cuadrática del número de neutrones presente en la sección eficaz de (CE$\nu$NS), esto analizando diferentes rangos de energía de retroceso nuclear. Para esto se emplearon flujos de neutrinos procedentes tanto de $\pi$-DAR (Decaimiento de pion en reposo) como de reactores. Por último, analizamos el efecto de las interacciones no estándar y la sensibilidad al ángulo de mezcla débil como ejemplos de las capacidades de estos futuros detectores.
